% Generated from locale/tr/content/normal-modal-logic/syntax-and-semantics/normal-modal-logics.tex; do not hand-edit.


\olfileid[tr]{nml}{syn}{nor}

\olsection{Normal Kiplik Mantıkları}

Bir kiplik mantığı, uygun biçimde iyi davranan herhangi bir !!{formula}s
kümesi olarak görülebilir. Hem ilginç yorumları ve uygulamaları bulunduğu,
fakat esas olarak kuramsal bakımdan iyi anlaşıldıkları için özellikle önem
taşıyan bazı kiplik mantıkları vardır. Bunlara normal kiplik mantıkları denir.

\begin{defn}\ollabel{defn:modal-logic}
Bir \emph{kiplik mantığı}~$\Log L$, temel kiplik dilinin aşağıdaki koşulları
sağlayan bir !!{formula}s kümesidir:
\begin{enumerate}
\item Bütün önermesel totolojileri içerir,
\item düzgün yerine koyma altında kapalıdır ve
\item modus ponens altında kapalıdır; yani $!A$ ve $(!A \lif
  !B) \in \Log L$ olduğunda $!B \in \Log L$ de olur.
\end{enumerate}
\end{defn}

\begin{ex}
Pek ilginç olmadıkları kabul edilse de bu tanımı sağlayan ve dolayısıyla
kiplik mantığı sayılan bazı !!{formula}s kümeleri şunlardır:
\begin{enumerate}
\item bütün totolojilerin kümesi,
\item bütün !!{formula}s kümesi (tutarsız kiplik mantığı).
\end{enumerate}
\end{ex}

\begin{defn}\label{defn:normal-modal-logic}
Bir kiplik mantığı~$\Log L$ ancak ve ancak aşağıdaki koşulları sağlıyorsa
\emph{normal}dir:
\begin{enumerate}
\item Aşağıdaki iki şemayı da içerir:
\begin{align}
\tag{K} \Box(p \lif q) & \lif (\Box p \lif \Box q) \\
\tag{Dual} \Diamond p & \liff \lnot \Box \lnot p
\end{align}
\item Zorunlulaştırma altında kapalıdır; yani $!A \in
  \Log L$ olduğunda $\Box !A \in \Log L$ de olur.
\end{enumerate}
\end{defn}

Kiplik mantıklarını tanımlamanın iki temel yoluyla karşılaşacağız. Bunlardan
biri, belirli !!{derivation} sistemlerinde türetilebilen !!{formula}s kümesi
olarak, ispat kuramsal yoldur. Diğeri ise belirli model sınıflarında geçerli
olan !!{formula}s kümeleri olarak, model kuramsal yoldur. Bu, olağan önerme
mantığındaki (hatta birinci dereceden mantıktaki) durumu yansıtır. Burada da
bir !!{formula}s kümesini iki ayrı yoldan belirleyebiliriz: örneğin bütün
doğruluk-değeri atamalarında doğru olan totolojilerin kümesi ve bir
!!{derivation} sistemindeki bütün !!{derivable} !!{formula}s kümesi olarak.
Normal kiplik mantıklarında bu iki yönlü belirleme, bağıntısal modeller ve
aksiyomatik (Hilbert tipi) ispat sistemleri kullanılarak yapılabilir.
